\section{The root cover and its Hodge line}
\label{sec:rankone}\label{hd:section}

We retain the diagram and the boundary conventions of
Proposition~\ref{it:diagram}.  Thus
\[
 X'\xrightarrow{x}W,\qquad h:W\longrightarrow Z,
 \qquad q=hx,
\]
and $X'$ dominates a fixed earlier log Iitaka model $X_*$.  Its boundary
$D'$ is obtained from the boundary on $X_*$ by strict transform and addition
of the reduced exceptional divisor.  Put
$D_W^0=(g^*D_Y)_{\mathrm{red}}$.  Lemma~\ref{it:zerofiber} gives, for very
general $z\in Z$ and for the geometric generic fiber $J$ of $x$,
\begin{equation}
 \begin{split}
 \kappa(X_{*,z},K_{X_{*,z}}+D_*|_{X_{*,z}})&=0,\\
 \kappa(W_z,K_{W_z}+D_W^0|_{W_z})&=0,\\
 \kappa(J)=\kappa(J,K_J+D'|_J)&=0.
 \end{split}
 \label{ro:initial-zero}
\end{equation}
All subsequent modifications are smooth projective birational
modifications, with this same source-boundary convention.  We shall
construct an effective rational boundary $B$ and a nef rational Hodge line
$M$ on $W$, prove a section comparison on $W_z$, and then show that the
Hodge line is rationally trivial there.

\subsection{An ordinary cyclic cover}

\begin{lemma}[The root cover]
\label{ro:root}
Let $p$ be the least positive integer for which
$H^0(J,pK_J)\ne0$.  A generator $\omega_J$ determines a geometrically
connected cyclic cover of $J$.  Every smooth projective resolution
$\widetilde J$ of its finite normalization satisfies
\[
 \kappa(\widetilde J)=0,
 \qquad h^0(\widetilde J,mK_{\widetilde J})=1
 \quad\text{for every }m>0.
\]
The root form spans the entire ordinary top-form space.  After shrinking
a nonempty open $W^\circ\subset W$, these covers admit an equivariant
resolution forming a smooth projective family over $W^\circ$.
\end{lemma}

\begin{proof}
The spaces of pluricanonical sections commute with extension to an
algebraic closure of the ground field.  We first make that extension and
work with the smooth projective variety $J$ over an algebraically closed
field of characteristic zero.  Since $\kappa(J)=0$, every nonzero
pluricanonical system has dimension one.

Choose a rational canonical frame $\tau$ and write
$\omega_J=f\tau^p$.  Normalize $J$ in the field obtained by adjoining a
$p$th root of $f$.  This field extension has degree $p$.  Indeed, since the
ground field contains the $p$th roots of unity, reducibility of the Kummer
polynomial would make $f$ a $q$th power for some prime $q\mid p$.  If
$f=a^q$, then the rational $(p/q)$-canonical form
$a\tau^{p/q}$ has regular $q$th power.  Its order at every prime is
therefore nonnegative, so it is regular, contrary to the minimality of
$p$.  This also proves geometric connectedness before the extension of
the ground field.

Denote the finite normalization by $\pi_0:J^{\mathrm{nor}}\to J$ and a
smooth resolution by $\mu:\widetilde J\to J^{\mathrm{nor}}$.  The rational
root form $\eta$ on $\widetilde J$ satisfies
\begin{equation}
 \eta^{\otimes p}=(\pi_0\mu)^*\omega_J
 \label{ro:root-identity}
\end{equation}
as pluricanonical differentials.  Pullback of a regular pluricanonical
differential under a morphism of smooth varieties is regular.  Thus
\eqref{ro:root-identity} proves that $\eta$ is regular at every prime of
$\widetilde J$, including the resolution-exceptional primes.

Set $E=\operatorname{div}(\omega_J)$.  At a prime of $J^{\mathrm{nor}}$
over a prime of $J$ at which $E$ has coefficient $a\ge0$, let $e$ be the
ramification index.  The tame differential formula gives
\begin{equation}
 \ord(\eta)=\frac{ea}{p}+e-1.
 \label{ro:root-order}
\end{equation}
If $a=0$, adjoining the root of a unit is unramified at that prime and
$e=1$.  If $a\ge1$, the right side of \eqref{ro:root-order} is at most
$2ea$.  This bound is a statement at the primes of the finite
normalization; we do not assert it at every prime of its resolution.

For $s\in H^0(\widetilde J,mK_{\widetilde J})$, the rational function
$s/\eta^m$ consequently has poles at the primes of $J^{\mathrm{nor}}$
bounded by $2m\pi_0^*E$.  Normality therefore gives an injection
\begin{equation}
 H^0(\widetilde J,mK_{\widetilde J})
 \longrightarrow
 H^0(J^{\mathrm{nor}},\cO_{J^{\mathrm{nor}}}(2m\pi_0^*E)),
 \qquad s\longmapsto s/\eta^m.
 \label{ro:root-section-injection}
\end{equation}
Only codimension-one tests on $J^{\mathrm{nor}}$ enter this injection.

Lemma~\ref{num:finite} gives
$\kappa(J^{\mathrm{nor}},\pi_0^*E)=\kappa(J,E)=0$.
Since $\pi_0^*E$ is effective, each space on the right of
\eqref{ro:root-section-injection} has dimension one: two independent
sections would give a nonconstant section ratio. The nonzero section
$\eta^m$ gives the reverse inequality on the left, proving the assertion.

The integer $p$ and the one-dimensional generator space are unchanged
geometrically.  A rational generator can consequently be chosen as a
section of $pK_{X'/W}$.  The Kummer construction is algebraic over a
nonempty open of $W$.  An equivariant resolution of its projective
normalization, followed by shrinking the base and generic smoothness,
gives the asserted smooth projective family, with its fiberwise cyclic
deck action.
\end{proof}


\subsection{The entire top line and its extension}

A highest Hodge line means the
highest nonzero filtration piece $F^p$ of a pure variation, with
$F^{p+1}=0$ and $\rank F^p=1$. The rank condition concerns the entire
highest step; a rank-one eigenspace in a larger highest step does not
satisfy it. Its projective period map is the map on the connected universal cover
to the projective space of a flat reference vector space determined
by that line.
Its differential rank need not equal the rank of the full period map.
Our applications use rational polarizations and a monodromy-invariant
lattice in the rational local system. The full-period comparison in Theorem~\ref{period:adjoint} also
allows a real polarization, with a lattice in the underlying real local
system. The parabolic extension is the rational line obtained by making
boundary monodromy unipotent on local root covers, taking its canonical
Hodge extension, and descending with the resulting rational weights.
Finite-monodromy weights are retained in this convention.

\begin{lemma}[Parabolic extension]\label{hodge:extensions}
Let $P$ be smooth projective, let $D$ be a reduced SNC divisor, and let
$L^\circ$ be the highest line of a polarized integral pure variation
on $P\setminus D$. Its rational parabolic extension $L$ is nef.
This extension commutes with positive tensor powers and with pullback
by morphisms of smooth projective log pairs whose interiors map into
$P\setminus D$.
\end{lemma}

\begin{proof}
For unipotent local monodromy, the extension is the top filtration
subbundle of the Deligne--Schmid extension. Nefness and compatibility
with powers and log-pair pullbacks are
\cite[Lemmas~2.16 and~5.6]{BFMT}.
The pullback statement follows from functoriality of the unipotent
extension and the extended Hodge filtration.

Integral polarized monodromy is quasi-unipotent at the boundary.
In a crossing chart, take sufficiently divisible roots of the boundary
coordinates to make it unipotent. The resulting line is equivariant
for the finite deck group. A common positive power kills the
stabilizer characters on its fibers and descends to a line bundle in
the original chart. These descents agree on overlaps: compare them on
a common root cover and use uniqueness of the unipotent extension.
They define the rational parabolic line, including its finite-monodromy
weights. The same comparison proves compatibility with powers and
further log-pair pullbacks. This proves the quasi-unipotent assertions
by descent from the unipotent ones.

These comparisons can be made on a global smooth adapted alteration.
For example, a finite neat level of the integral monodromy, followed
by compactification and resolution, makes the boundary monodromies
unipotent. Extra boundary components on which the original variation
already extends have weight zero. The pullback of $L$ is the unipotent
extension there and is nef. Nefness descends under a proper surjective
map: every curve downstairs is dominated by a curve upstairs, and the
projection formula compares their degrees. Thus $L$ is nef.
\end{proof}

Fix the rational relative form $\omega_J$ of Lemma~\ref{ro:root}.  Enlarge
an SNC divisor on $W$ so that the root-cover family and its pure integral
polarized middle-cohomology variation are defined on its complement.
This enlargement specifies the domain of the variation; it does not
add logarithmic poles to the target section systems.

Let $\mathcal L$ be the highest nonzero Hodge filtration piece of the
middle-cohomology variation of the resolved root-cover family.
Lemma~\ref{ro:root} says that $\mathcal L$ has rank one.  Write $M$ for
its parabolic rational extension to the prepared SNC compactification
$W$.  Lemma~\ref{hodge:extensions} makes $M$ nef.

\subsection{Actual divisorial orders and the boundary}

Two divisorial tests enter the comparison. The first records which
poles a coefficient on $W$ may have while its product with the
relative pluriform remains logarithmic at every actual source
component. It gives a minimum over those components. The second
test ranges over divisorial valuations on all higher models; the
metric calculation below will identify it with an integrability
threshold. For each prime $P$ of $W$, we denote these two quantities
by $t_P$ and $\lambda_P$, respectively.

\begin{proposition}[Divisorial normalization]
\label{ro:divisors}
For a prime divisor $P$ on $W$, and for a prime divisor $Q$ on the actual
source $X'$ dominating $P$, set
\begin{equation}
 \begin{aligned}
 m_Q&=\ord_Q(x^*P),\qquad
 r_Q=\frac1p\ord_Q(\omega_J),\qquad
 d_Q=\operatorname{coeff}_Q(D'),\\
 t_P&=\min_{\substack{Q\mapsto P\\Q\text{ on }X'}}
                   \frac{r_Q+d_Q}{m_Q},\\
 \lambda_P&=\inf_{Q\mapsto P}\frac{1+r_Q}{m_Q},\qquad
 B_P=1-\lambda_P+t_P.
 \end{aligned}
 \label{ro:orders}
\end{equation}
In the definition of $\lambda_P$, $Q$ runs over divisorial valuations on all higher
smooth source models.  In both formulas $r_Q$ is measured in the actual
relative canonical bundle with respect to the smooth base $W$.

The two rational divisors
\[
 T=\sum_Pt_PP,
 \qquad B=\sum_PB_PP
\]
have finite support.  The divisor $B$ is an effective boundary and
\begin{equation}
 0\le B_P\le1,
 \qquad B\ge D_W^0=(g^*D_Y)_{\mathrm{red}}.
 \label{ro:boundary}
\end{equation}
The formulas at unchanged codimension-one base points are invariant
under further source resolutions with the prescribed boundary
convention.  They are also preserved, allowing splitting of components,
by an \'etale extension of the base DVR and the corresponding source
base change.  After a further birational preparation of $W$, $B$ can be
assumed to have SNC support.
\end{proposition}

\begin{proof}
At the generic point of $P$, correct the relative form by the factor
$u^{-pt_P}$, where $u$ is a local equation for $P$; take a tensor power to
make the exponent integral.  At every actual component over $P$ its
normalized order satisfies
\begin{equation}
 r_Q-m_Qt_P+d_Q\ge0.
 \label{ro:actual-test}
\end{equation}
Orders at divisors horizontal over $W$ are nonnegative, since
$\omega_J$ is an ordinary regular pluricanonical form on the generic
fiber.  Near the generic point of $P$, the corrected form is therefore a
logarithmic pluriform for the actual SNC pair $(X',D')$.

The discrepancy inequality for a reduced SNC pair says that pullback
to a higher smooth model is still logarithmic for strict transform plus
reduced exceptional boundary.  In particular, the normalized ordinary
order of the corrected form at any higher prime is at least $-1$.
It follows that
\[
 \frac{1+r_Q}{m_Q}\ge t_P
 \quad\text{for every }Q\mapsto P,
 \qquad\text{and hence}\qquad \lambda_P\ge t_P.
\]
If $Q$ attains the actual minimum defining $t_P$, then
\begin{equation}
 0\le\lambda_P-t_P
 \le \frac{1+r_Q}{m_Q}-\frac{r_Q+d_Q}{m_Q}
 =\frac{1-d_Q}{m_Q}\le1.
 \label{ro:boundary-inequalities}
\end{equation}
This proves $0\le B_P\le1$.  If $P$ is contained in $g^{-1}D_Y$, every
actual component over $P$ lies in $D'$ and has $d_Q=1$.  The middle term
of \eqref{ro:boundary-inequalities} is then zero, so $B_P=1$.

Resolve, over the generic point of $P$, the support of the form's zero
and pole divisor together with $x^*P$.  On this resolution the relevant
conditions are monomial inequalities.  The SNC discrepancy inequality
shows that inequalities at its components imply the inequalities on
every higher model.  Thus a finite log resolution computes $\lambda_P$;
in particular it is rational.  Outside a finite divisorial subset of
$W$, the fibers are smooth and integral, the relative form has no
vertical zero or pole, and there is no vertical source boundary.  There
the actual component has $m_Q=1$, $r_Q=d_Q=0$, and the smooth-fiber
calculation gives $t_P=0$ and $\lambda_P=1$.  This proves finite support.

On a further source resolution, the old actual divisors remain.  They
still attain the old minimum.  Equation~\eqref{ro:actual-test} holds at
every new actual divisor by logarithmic pullback, so the minimum cannot
decrease.  The all-valuation infimum is unchanged.  An \'etale base-DVR
extension preserves canonical frames up to units, multiplicities, and
all the displayed orders; the same assertion holds on a fixed log
resolution, including when a divisor splits.  Finally resolve the finite
support on the base, resolve the source map, and recompute
\eqref{ro:orders}.  All previously unchanged codimension-one points keep
their coefficients, and new support lies on the exceptional base
divisor.  The inequalities just proved apply to the recomputed data.
\end{proof}

\begin{corollary}[A zero boundary coefficient]
\label{ro:zero-boundary}
If an actual coefficient $B_P$ in Proposition~\ref{ro:divisors} is zero,
some actual source component $Q$ over $P$ satisfies
\begin{equation}
 d_Q=0,\qquad m_Q=1,\qquad
 \frac{1+r_Q}{m_Q}
   =\inf_{Q'\mapsto P}\frac{1+r_{Q'}}{m_{Q'}}.
 \label{ro:zero-minimizer}
\end{equation}
The same conclusion applies after the \'etale DVR extensions described
there.  It refers to the actual coefficient before any operation that
assigns a new reduced-exceptional boundary coefficient.
\end{corollary}

\begin{proof}
Take $Q$ attaining the first minimum in \eqref{ro:orders}.  The equality
$B_P=0$ makes $\lambda_P-t_P$ in
\eqref{ro:boundary-inequalities} equal to one.  Thus
\[
 1\le\frac{1-d_Q}{m_Q}\le1.
\]
Since $m_Q$ is a positive integer and $d_Q\in\{0,1\}$, this forces
$m_Q=1$ and $d_Q=0$.  Equality throughout also gives the final assertion
in \eqref{ro:zero-minimizer}.
\end{proof}

The component singled out by Corollary~\ref{ro:zero-boundary} will
enter the ramification argument in Sections~\ref{sec:marking}--\ref{ds:section}.
That argument uses both its multiplicity one and its attainment of
the all-valuation threshold. We first use the two order tests to
identify the Hodge line and compare sections.

\subsection{The metric identifies the Hodge contribution}

\begin{proposition}[The Hodge order]
\label{ro:metric}
At every prime $P$ of $W$, the parabolic order of the root form $\eta$ is
$\lambda_P-1$.  Consequently
\begin{equation}
 M\sim_{\Q}\sum_P(\lambda_P-1)P,
 \qquad
 T\sim_{\Q}B+M.
 \label{ro:metric-identity}
\end{equation}
\end{proposition}

\begin{proof}
On the smooth family locus, finite change of variables gives
\begin{equation}
 \|\eta_w\|_{\mathrm{Hdg}}^2
   =c_p\int_{J_w}|\omega_{J,w}|^{2/p},
 \label{ro:norm-fiber}
\end{equation}
where $c_p>0$ is constant and the expression on the right is the
canonical density associated with a pluricanonical form.  In a local
canonical frame, this density is the $2/p$ power of its coefficient
times the usual canonical volume density.

Take a general point of $P$, with base coordinates $(u,s_2,\ldots,s_l)$
and $P=(u=0)$.  Integrate \eqref{ro:norm-fiber} against base Lebesgue
measure and the factor $|u|^{-2\gamma}$.  Fubini and change of variables
reduce local integrability to an integral on a log resolution of the
total space.  The factors from the absolute canonical density and from
the base canonical frame cancel precisely as prescribed by
$K_{X'/W}=K_{X'}-x^*K_W$.  Thus, at a vertical prime $Q$, the monomial
exponent in the resulting density is
\[
 r_Q-\gamma m_Q.
\]
The transverse integral is finite exactly when
$r_Q-\gamma m_Q>-1$.  Horizontal primes have nonnegative ordinary order
and cause no further condition.  A finite SNC resolution computes all
these conditions, so the supremal integrability exponent is
\begin{equation}
 \sup\left\{\gamma:
 |u|^{-2\gamma}\|\eta\|_{\mathrm{Hdg}}^2\in L^1_{\mathrm{loc}}
 \right\}
 =\inf_{Q\mapsto P}\frac{1+r_Q}{m_Q}
 =\lambda_P.
 \label{ro:threshold}
\end{equation}
Here local integrability is measured against base Lebesgue measure.
The value at the endpoint itself is irrelevant to this supremum.

We compare this with the parabolic extension.  After a transverse root
coordinate killing the finite part of the monodromy, a nonvanishing
canonical extension frame has squared Hodge norm bounded above and
below by powers of $-\log|u|$.  This is the one-variable
nilpotent-orbit growth estimate for a polarized pure variation
\cite[\S\S4 and 6]{Schmid}.  Accordingly, a rational section of
parabolic order $a$ has squared norm
\[
 |u|^{2a}\,(-\log|u|)^{O(1)},
\]
where the notation means two-sided bounds by fixed powers of the
logarithm.  Its supremal exponent in \eqref{ro:threshold} is $1+a$.

There is a meromorphic coefficient to which this statement applies.
The form $\eta$ is an algebraic section of the top line on the
family locus.  One may use the regular-singular extension of the
geometric Gauss--Manin system.  Equivalently, the finiteness of the
threshold in \eqref{ro:threshold} gives a weighted $L^2$ bound for its
holomorphic coefficient in a nonvanishing extension frame.  Absorb the
logarithmic factors into an arbitrarily small power of the radius.  The
Laurent-series integral test then bounds the negative exponents of that
coefficient, excluding an essential singularity and proving meromorphic
extension.  The finite root coordinate gives a rational parabolic order
downstairs.  Clearing its denominator by a tensor power, and applying
the codimension-one extension test, gives the divisor of the rational
line $M$.

Comparing with \eqref{ro:threshold} proves $a=\lambda_P-1$.
There is no additional factor of $p$: the normalized order $r_Q$ already
contains $1/p$, and the density in \eqref{ro:norm-fiber} uses the same
normalization.  Finally
$t_P=B_P+(\lambda_P-1)$ at every prime, which proves
\eqref{ro:metric-identity}.  The transverse root coordinate has only
been used to compute the metric order; no section system has been
descended from a ramified alteration.
\end{proof}

\subsection{Preparation and transfer of sections}

We first record the birational preparation needed to avoid an
uncontrolled codimension-two image in a section comparison.

\begin{lemma}[Extracting source divisorial valuations]
\label{ro:extraction}
Let $T_*$ be a fixed smooth projective model of the source of a dominant
rational map to a smooth projective variety $R$.  After resolving the
map, making a birational modification of $R$, and resolving the source
again, we can arrange that every source divisor whose image in the new
base has codimension at least two is exceptional over $T_*$.  Further
smooth birational preparations of the base and compatible source
resolutions preserve this assertion.

If the maps also commute with maps to a fixed base $Z$, the same
assertion holds on their very general fibers over $Z$.
\end{lemma}

\begin{proof}
Resolve the rational map, obtaining a morphism $T_1\to R$.  Put
$e=\dim T_1-\dim R$.  A prime divisor of $T_1$ with image of codimension
at least two lies in the closed locus where the fiber dimension is at
least $e+1$.  There are only finitely many divisorial components in that
locus.  Consider those not exceptional over $T_*$.

For each such prime, its divisorial valuation $v$ on
$K=\C(T_1)$ has nontrivial restriction $w$ to $L=\C(R)$.  The latter is
a discrete rank-one valuation.  We justify that it is divisorial.
Residue elements algebraically independent over $\kappa(w)$ lift to
elements algebraically independent over $L$: in a proposed relation,
divide by a coefficient of least valuation and reduce to the residue
field.  Hence
\[
 \trdeg_{\kappa(w)}\kappa(v)\le\trdeg_LK.
\]
Since $v$ is divisorial, this gives
\[
 \trdeg_{\C}\kappa(w)
 \ge (\dim T_1-1)-(\dim T_1-\dim R)=\dim R-1.
\]
Conversely, lifts of algebraically independent residue elements for $w$,
together with an element of positive valuation, are algebraically
independent over $\C$.  Thus
$\trdeg_{\C}\kappa(w)\le\dim R-1$.  Equality follows.

Choose a residue transcendence basis and lift it to units of $L$,
together with an element generating the value group of $w$.  They give
a purely transcendental subfield of $L$ over which $L$ is finite.
Normalize the corresponding model in $L$ and take the center of $w$
above the DVR cut out by that positive-valuation parameter.  This
realizes $w$ as a prime divisor on a birational model of $R$.
Extract all the finitely many restrictions in this way and dominate the
result by a smooth projective model.  On a compatible smooth source
model, each of the original nonexceptional divisors now has divisorial
image.  Any new source prime is exceptional over $T_*$.  A further
birational base modification preserves the strict transform of each
extracted valuation divisor, proving persistence.

For the last assertion, shrink a nonempty open in $Z$ so that the
finitely many relevant loci, their images, and their boundary strata
have the expected fiber dimensions.  A prime divisor contracted to
codimension at least two on a fiber lies in the restriction of the
fixed excess-fiber-dimension locus.  A global component of codimension
at least two dominating $Z$ still has codimension at least two on a
general fiber.  Thus a divisorial component on a general source fiber
comes from a global divisorial component of that locus, possibly after
splitting.  Its image over $T_*$ has codimension at least two, and the
same dimension calculation makes its image exceptional on the general
fiber of $T_*$.  This proves the fiberwise assertion.
\end{proof}

\begin{proposition}[The Iitaka-fiber section bound]
\label{ro:transfer}
The models in Proposition~\ref{it:diagram} can be chosen, with the
divisors of Proposition~\ref{ro:divisors} recomputed on them, so that
\begin{equation}
 \kappa\bigl(W_z,K_{W_z}+B|_{W_z}+M|_{W_z}\bigr)\le0
 \label{ro:transfer-bound}
\end{equation}
for very general $z\in Z$.  The comparison is a comparison with the
section system on the fixed fiber $X_{*,z}$ in
\eqref{ro:initial-zero}.
\end{proposition}

\begin{proof}
Apply Lemma~\ref{ro:extraction} with fixed source $X_*$ and base $W$.
Retain the maps to $Y$ and $Z$, resolve the finite boundary supports,
and recompute \eqref{ro:orders}.  By the lemma, every divisor of a very
general fiber $X'_z$ whose image under $x_z$ has codimension at least two
is exceptional over $X_{*,z}$.

We also shrink the good open in $Z$ for the finitely many supports and
resolutions used in the divisorial calculation.  The tests defining $T$
then restrict to $W_z$.  To see this explicitly, over the smooth good
locus of $W$ those tests are zero.  Outside it, only restrictions of the
finitely many divisors dominating $Z$ can meet a general $W_z$ in a
divisor.  Generic smoothness and adjunction preserve the multiplicities
and relative-canonical orders there, with splitting of components
allowed.  Loci of higher codimension stay of higher codimension after
the shrinking just made.  The same reasoning on a fixed log resolution
computes all the higher-valuation inequalities.

Write $T_z=T|_{W_z}$.  For a sufficiently divisible integer $m$, take
\[
 \sigma\in H^0(W_z,m(K_{W_z}+T_z)).
\]
Using the relative canonical identification, multiply its pullback by
$\omega_J^{m/p}$.  This is a rational $m$-canonical form on $X'_z$.
If $Q$ is an actual source divisor with divisor image $P$, its normalized
order, after allowing the source boundary, is at least
\[
 m\bigl(r_Q-m_Qt_P+d_Q\bigr)\ge0
\]
by \eqref{ro:actual-test}.  A horizontal divisor passes the ordinary
test since the generic fiber generator is regular.  All other
untested divisors are exceptional over $X_{*,z}$.

Push the rational form to $X_{*,z}$.  At every prime of that smooth
fixed fiber, its strict transform has just passed the required
logarithmic test; hence the pushdown belongs to
\[
 H^0(X_{*,z},m(K_{X_{*,z}}+D_*|_{X_{*,z}})).
\]
Possible poles on a divisor exceptional over $X_{*,z}$ have no bearing
on this codimension-one test downstairs.  The construction is
injective.  More precisely, the common factor $\omega_J^{m/p}$ cancels
in ratios, so ratios of two input sections pull back by the dominant
map $x_z$ and their rational-map dimension is preserved.  The target
system has Iitaka dimension zero by \eqref{ro:initial-zero}.  The
identity $T_z\sim_{\Q}B|_{W_z}+M|_{W_z}$ from
Proposition~\ref{ro:metric} proves \eqref{ro:transfer-bound}.
\end{proof}

\subsection{Adjoint positivity for the restricted variation}

The section bound is now available on a very general $h$-fiber.
On that fiber, $K+B$ has a nonzero divisible section because
$B\geq D_W^0$. We shall use adjoint positivity to show that any
positive-dimensional adjoint base for the restricted Hodge line
would contradict the bound for $K+B+M$. The following is the
full companion input; its projective form follows immediately after.

\begin{theorem}[Full-period adjoint comparison]\label{period:adjoint}
Let $B$ be a smooth compact connected manifold in Fujiki class
$\mathcal C$, and let $B^\circ\subset B$ be a dense Zariski open
whose complement is an SNC divisor. Let $\mathbb V$ be a
real-polarizable pure variation of Hodge structures on $B^\circ$
with an integral lattice. Suppose a complex direct summand of
$\mathbb V_{\C}$ has a highest nonzero Hodge filtration step of rank
one, with rational parabolic extension $L$. There exist a modification
$\tau:B'\to B$, a connected-fiber surjective morphism $p:B'\to R$ to
a smooth projective variety, and a nef rational line bundle $A$ on
$R$ such that
\[
 \tau^*L=p^*A\quad\text{in }\operatorname{Pic}(B')\otimes\Q,
 \qquad K_R+cA\text{ is big for some }c\in\Q_{>0}.
\]
Positive rational rescaling of $L$ is allowed, and the line identity
persists under further modifications. A point is allowed for $R$;
then $\tau^*L$ is rationally trivial.
\end{theorem}

\begin{proof}
This is \cite[\PoneAdjoint]{SubadditivityCompanion},
whose proof uses the full period image and boundary numerical-rank
loss. The statement permits a complex summand of an integral ambient
variation; it is stronger in this respect than a theorem requiring
that summand itself to carry an integral lattice.
\end{proof}

\begin{proposition}[Projective adjoint comparison]\label{hd:adjoint}
Let $P$ be smooth connected projective, let $D$ be reduced SNC,
and let $L$ be the rational parabolic extension of the entire
highest Hodge line of a polarizable integral pure variation on
$P\setminus D$. There are a birational morphism
$\mu:P'\to P$ with $P'$ smooth projective, a surjective morphism
$a:P'\to R$ with connected fibers to a smooth projective variety,
and a nef rational line bundle $A$ on $R$ such that
\begin{equation}\label{hd:adjoint-conclusion}
 \mu^*L\sim_{\mathbf Q}a^*A,
 \qquad K_R+bA\text{ is big for all sufficiently large rational }b.
\end{equation}
The point quotient is allowed, in which case $\mu^*L$ is
rationally trivial.
\end{proposition}

\begin{proof}
The hypotheses meet Theorem~\ref{period:adjoint} directly. A projective
manifold belongs to class $\mathcal C$, a rational polarization
is a real polarization, and the entire complexification of the
given variation is a complex direct summand of itself. Its
highest nonzero filtration step is the specified rank-one line.
Thus the companion supplies $\tau:B'\to P$, $p:B'\to R$,
$A$ and $c>0$ with
\begin{equation}\label{direct:identity}
                \tau^*L\sim_{\mathbf Q}p^*A,
                \qquad K_R+cA\text{ big}.
\end{equation}

It remains to obtain a projective source model without changing
this identity or the connected-fiber property. The map
$P\dashrightarrow R$ induced by $p$ is meromorphic between
projective varieties and hence rational: its graph closure is
algebraic. Resolve that graph projectively to obtain a smooth
projective $P'$ and morphisms
\[
                       \mu:P'\to P,
                       \qquad a:P'\to R.
\]
On a common smooth modification of $B'$ and $P'$, the two maps
to $R$ agree. Pulling back \eqref{direct:identity} there therefore
identifies the pullbacks of $\mu^*L$ and $a^*A$.

Pullback by a proper birational morphism $v:T\to P'$ is
injective on rational line bundles. Indeed, if a line bundle
$N$ pulls back to a torsion line bundle, some positive power
$v^*N^m$ is trivial. Normality of $P'$ gives
$v_*\mathcal O_T=\mathcal O_{P'}$, and the projection formula
then gives $N^m\simeq\mathcal O_{P'}$. Applying this observation
to the common modification proves
\[
                          \mu^*L\sim_{\mathbf Q}a^*A.
\]

The generic fibers of $a$ are connected as well. The original
map $p$ has connected fibers and normal target, so
$\mathbf C(R)$ is relatively algebraically closed in
$\mathbf C(B')=\mathbf C(P)$. The same field extension occurs
for $a$. Its Stein factor is consequently finite birational
over the normal variety $R$, and hence is $R$ itself. Thus all
fibers of $a$ are connected.

Finally $A$ is nef. For every rational $b\geq c$,
$K_R+bA=(K_R+cA)+(b-c)A$ is big. If $R$ is a point, the line
identity already gives rational triviality; the zero-dimensional
bigness statement uses the usual point convention. This proves
the proposition.
\end{proof}

\begin{remark}\label{hd:further-models}
The conclusion is stable under further smooth projective
birational modification $\pi:R'\to R$ and a compatible
modification of the source.  The parabolic pullback property
preserves the line-bundle identity and nefness, while
\[
 K_{R'}+b\pi^*A
       =\pi^*(K_R+bA)+E_\pi,\qquad E_\pi\geq0,
\]
preserves bigness.  This is the form used when preparing the
later section comparisons.
\end{remark}

\subsection{Rational triviality on the Iitaka fiber}

The following weak-positivity consequence specifies the section
arithmetic used in the argument.

\begin{lemma}[Paying for a weak-positivity twist]
\label{ro:weak-addition}
Let $a:T_1\to R$ be a surjective morphism between smooth connected
projective varieties.  Let $\Delta$ be an effective rational SNC
boundary with coefficients in $[0,1]$, and let $N\sim_{\Q}a^*A$ for a
nef rational divisor $A$ on $R$.  Suppose
\[
 \kappa(T_1,K_{T_1}+\Delta)\ge0,
 \qquad K_R+cA\text{ is big}
\]
for a positive rational number $c$.  Then
\begin{equation}
 \kappa(T_1,K_{T_1}+\Delta+cN)\ge\dim R.
 \label{ro:weak-addition-bound}
\end{equation}
\end{lemma}

\begin{proof}
Use Lemma~\ref{ro:extraction} with fixed source $T_1$ and base $R$.
On the resulting smooth source $T_2$, let $\Delta_2$ be the strict
transform of $\Delta$ plus the reduced exceptional divisor, and pull
back $N$ and $A$; write $N_2=\pi^*N$ for the former pullback.
The SNC discrepancy inequality gives
\[
 K_{T_2}+\Delta_2
   =\pi^*(K_{T_1}+\Delta)+E_{\mathrm{exc}},
 \qquad E_{\mathrm{exc}}\ge0\text{ exceptional over }T_1.
\]
The section systems in the assertion therefore have the usual
pullback and pushdown comparison.  For a smooth birational base map
$\rho:R'\to R$,
\[
 K_{R'}+c\rho^*A
   =\rho^*(K_R+cA)+K_{R'/R}
\]
is still big, because $K_{R'/R}$ is effective.  Relabel the prepared
base as $R$, its divisor as $A$, and its source morphism as $a_2$.

Choose a sufficiently divisible positive integer $m$ for which the
relative direct image
\[
 \mathcal E_m
   =(a_2)_*\cO_{T_2}
          \bigl(m(K_{T_2/R}+\Delta_2)\bigr)
\]
has positive rank.  Such a choice is possible by restriction of a
nonzero section of a multiple of $K_{T_1}+\Delta$ to the generic fiber,
followed by pullback.  The pair $(T_2,\Delta_2)$ is log canonical,
$T_2$ is projective, and the base is smooth projective.  Thus
\cite[Theorem 1.1]{FujinoWeakPositivity} makes $\mathcal E_m$ weakly positive.

Put $L=K_R+cA$ and fix an ample integral divisor $H$ on $R$.  Increase
$m$ to clear the denominators needed below.  Since $L$ is big, choose a
positive integer $\alpha$ such that
\[
 \alpha mL-H\quad\text{is big}.
\]
The definition of weak positivity
\cite[Definition 7.2]{FujinoWeakPositivity} supplies a positive integer $\beta$ for
which
\begin{equation}
 \bigl(\operatorname{Sym}^{\alpha\beta}\mathcal E_m\bigr)^{**}
       \otimes\cO_R(\beta H)
 \quad\text{is generically generated by global sections}.
 \label{ro:weak-generation}
\end{equation}
At the generic point of $R$, multiplication of relative sections is
not the zero map: the power of any nonzero section on the integral
generic fiber is nonzero.  Generic generation in
\eqref{ro:weak-generation} therefore yields a global section whose
evaluation is a nonzero rational section $u$ of
\[
 \alpha\beta m(K_{T_2/R}+\Delta_2)+\beta a_2^*H.
\]
At every codimension-one point of $R$, the direct image is locally
free and the reflexive symmetric power agrees with its ordinary
symmetric power.  Evaluation is regular above those points.  Any
remaining pole divisor of $u$ thus has image of codimension at least
two and is exceptional over the fixed source $T_1$.

For a sufficiently divisible positive integer $q$, multiply $u^q$ by
the pullbacks of all sections of
\[
 q\beta(\alpha mL-H).
\]
The resulting rational sections lie in the class
\begin{align*}
 &q\alpha\beta m(K_{T_2/R}+\Delta_2)
       +q\beta a_2^*H
       +q\beta a_2^*(\alpha mL-H)\\
 &\hspace{25mm}
   =q\alpha\beta m(K_{T_2}+\Delta_2+cN_2).
\end{align*}
They are regular at every prime not exceptional over $T_1$, and so
push down to the required absolute section system on $T_1$.  Their
ratios are the pullbacks of the ratios of sections of the big divisor
$\alpha mL-H$.  These ratios give a rational map of dimension
$\dim R$, proving \eqref{ro:weak-addition-bound}.
\end{proof}

\begin{proposition}[Triviality on the Iitaka fiber]
\label{ro:trivial}
For very general $z\in Z$,
\begin{equation}
 M|_{W_z}\sim_{\Q}0,
 \qquad
 \kappa(W_z,K_{W_z}+B|_{W_z})=0.
 \label{ro:trivial-identities}
\end{equation}
These assertions also hold on the geometric generic $h$-fiber.
\end{proposition}

\begin{proof}
Set $T_1=W_z$, $\Delta=B|_{W_z}$, and $N=M|_{W_z}$. For the chosen
$z$, the fiber is smooth and all relevant boundary strata meet it
transversely, including those of the enlarged divisor defining the
domain of the variation. The log-pair pullback property in
Lemma~\ref{hodge:extensions} therefore identifies
$N$ with the parabolic extension of the restricted variation.
Moreover
\[
 \kappa(T_1,K_{T_1}+\Delta)\ge0
\]
by \eqref{ro:initial-zero} and $B\ge D_W^0$.  In particular,
$D=K_{T_1}+\Delta$ is rationally equivalent to an effective rational
divisor.

Apply Proposition~\ref{hd:adjoint} to the restricted integral polarized
variation whose \emph{entire} highest piece is this line. The comparison is applied directly to this restricted variation;
its entire complexification is an allowed summand in the companion theorem. On a higher smooth model it gives
\[
 N\sim_{\Q}a^*A,
 \qquad a:T_1\longrightarrow R,
 \qquad A\text{ nef},
 \qquad K_R+cA\text{ big for }c\gg1.
\]
On the higher source model use strict transform plus reduced
exceptional boundary for $\Delta$.  The logarithmic section comparison
preserves \eqref{ro:transfer-bound} and the rational effectivity of
$D$; the Hodge line is pulled back.  If $\dim R>0$, choose a sufficiently
large rational $c\ge1$.  Lemma~\ref{ro:weak-addition} gives
\[
 \kappa(T_1,D+cN)\ge\dim R>0.
\]
But
\begin{equation}
 D+N=\frac1c(D+cN)+\left(1-\frac1c\right)D.
 \label{ro:interpolation}
\end{equation}
After clearing denominators, multiplication by the section of the
effective last summand preserves rational-map dimensions.  Thus the
same positive lower bound holds for $D+N$, contradicting
\eqref{ro:transfer-bound}.  We conclude that $R$ is a point and
$N\sim_{\Q}0$.  Rational triviality descends through the intervening
birational morphisms, since their direct image of the structure sheaf
is the structure sheaf.  Now \eqref{ro:transfer-bound} and
$\kappa(D)\ge0$ give $\kappa(D)=0$.

Finally, these algebraic assertions pass to the geometric generic
fiber.  For the Kodaira dimension assertion, use generic base change
in the countably many divisible degrees.  For rational triviality,
choose an integer clearing the denominator of $M$.  For each positive
integer $j$, simultaneous nonvanishing of sections of the $j$th power
of that line and of its inverse is a closed condition on a smooth
proper family, and is equivalent to triviality on a connected
projective fiber.  The countable union of these loci contains every
very general $z$ by what we proved.  Hence some one of them contains a
dense open of $Z$; otherwise a very general point would avoid them
all.  Generic base change gives the required rational triviality on
the geometric generic fiber.
\end{proof}

\subsection{Flatness and finite character}

Rational triviality gives degree zero on every complete curve in an
Iitaka fiber. The curvature calculation below converts those zero
degrees into a flat highest line; finite character is then a statement
about the integral variation restricted to that algebraic locus.

\begin{corollary}[Degree zero on a curve]\label{hodge:degree-zero}
In Lemma~\ref{hodge:extensions}, suppose that $P$ is a smooth
projective curve. If $\deg L=0$, the projective highest-line map is
constant on the connected universal cover of $P\setminus D$.
Equivalently, $L^\circ$ is a flat line subbundle.
\end{corollary}

\begin{proof}
Pass to a finite cover making the boundary monodromies unipotent.
The extension pulls back, and its degree remains zero. Let $\alpha$
be the highest-line Hodge curvature. Griffiths' curvature formula
makes it semipositive and identifies its kernel with the kernel of
the projective line differential.

We check that its integral is the degree of the extension. For a smooth
reference curvature $\beta$, write $\alpha-\beta=\ddc u$. Near a
puncture, with $\rho=-\log|z|^2$, the Hodge norm and horizontal Schwarz
estimates give
\[
 |u|\le C(1+\log\rho),\qquad
 0\le\alpha\le C\frac{i\,dz\wedge d\bar z}{|z|^2\rho^2};
\]
see \cite[Theorem~5.1]{CK1989} and
\cite[Proposition~2.4 and Lemma~2.5]{BrunebarbeCadorel}.
Use products of cutoffs $\chi((\log\rho)/N)$, equal to one away
from a shrinking puncture neighborhood. Their second derivatives
are uniformly bounded in the cusp metric and supported on escaping
tails. Integration by parts bounds the error by the tail integral of
$(1+\log\rho)\rho^{-2}\,d\rho$, which tends to zero.
Thus $\int_{P\setminus D}\alpha=\deg L=0$. A continuous
nonnegative curvature form with zero integral vanishes. Its kernel
identity proves constancy on the cover and therefore on the original
universal cover.
\end{proof}

\begin{corollary}[Finite character of a line]\label{lem:finite-character}
Let $\mathbb V$ be a polarized integral variation of Hodge structures
on a connected smooth complex algebraic variety. Every rank-one
complex local subsystem of $\mathbb V_{\C}$ has finite monodromy.
\end{corollary}

\begin{proof}
Deligne's determinant theorem
\cite[Corollary~4.2.8(iii)(b)]{Deligne71} states that a positive tensor
power of the determinant of any complex local subsystem is trivial.
For a rank-one subsystem, the determinant is the subsystem itself.
Its character therefore takes values in a fixed finite group of roots
of unity. The theorem applies anew to the pulled-back variation on
any connected smooth algebraic locus.
\end{proof}

\begin{corollary}[Restriction and finite character]
\label{hodge:restricted-character}
Let $T$ be a connected smooth complex algebraic variety mapping into
the domain of a polarized integral pure variation whose entire highest
Hodge step is a line. If the pulled-back projective highest-line map is
constant, its highest line has finite monodromy character.
\end{corollary}
\begin{proof}
The pulled-back variation is again polarized, integral and pure.
Constancy of its projective highest line makes that line a rank-one
complex local subsystem. Apply Corollary~\ref{lem:finite-character}
to this restricted variation. No common exponent for different
algebraic loci is needed.
\end{proof}

\begin{corollary}[Flatness and finite character on the Iitaka fiber]
\label{ro:flat-character}
For very general $z\in Z$, the projective entire top line of the
restricted root-cover variation is constant on the connected good
locus of $W_z$, and its top-line character has finite image.
\end{corollary}

\begin{proof}
Put $T_1=W_z$ and $N=M|_{W_z}$. Proposition~\ref{ro:trivial}
gives $N\sim_{\Q}0$, and the pullback property in
Lemma~\ref{hodge:extensions} identifies $N$ with the parabolic
extension of the restricted entire top line.
The degree of $N$ on every complete smooth test curve is zero.
Corollary~\ref{hodge:degree-zero} makes the projective top line constant
along such curves.  Complete-intersection curves through general
points and tangent directions give projective constancy on the
connected good locus of $T_1$.
Corollary~\ref{hodge:restricted-character}, applied to this
restricted integral polarized variation, gives finite monodromy for
its original top line.
\end{proof}

\subsection{Alternative arguments for the Hodge line}
\label{ro:unitary-character}

\begin{remark}[Tensor comparison on restrictions]
The direct comparison above was applied to the restricted variation.
A separate integral tensor replacement, which makes
projective monodromy faithful, is proved in
\cite[\PfourTensor]{ProjectiveMethodsCompanion}; it gives the independent
adjoint comparison in \cite[\PfourAdjoint]{ProjectiveMethodsCompanion}.
On a connected algebraic restriction, that tensor construction must
be applied anew to the pulled-back integral variation. A replacement
faithful for the ambient moving line need not remain faithful after
restriction, so the monodromy and rational isotypic calculation must
be repeated on the restricted locus. Its positive tensor exponent may
depend on that locus; no common exponent is required. Thus it gives
an alternative adjoint comparison on each restriction. The proof here
uses the direct comparison and Corollary~\ref{hodge:restricted-character}.
\end{remark}

There is also a direct metric proof of the finite-character conclusion
in Corollary~\ref{ro:flat-character}. It uses the rational triviality
from Proposition~\ref{ro:trivial}, in addition to zero curvature.
Put $T_1=W_z$ and $N=M|_{W_z}$ on the prepared smooth
projective model. The constant projective top line on its connected
good open is a flat line. Its polarization, with the fixed Hodge sign,
is a positive flat metric, so its monodromy character $\chi$ is
unitary. Quasi-unipotence of the integral ambient variation makes
the characters around the finitely many SNC boundary components
roots of unity \cite{Schmid}.

Choose $e>0$ killing these local characters and clearing the
parabolic weights. Then $\chi^e$ extends across the boundary to
a unitary flat line on $T_1$; its holomorphic line bundle is the
parabolic tensor power $eN$. Since $N\sim_{\Q}0$, a further
positive tensor power is holomorphically trivial. For a nowhere-zero
holomorphic section of that power, the logarithm of its norm is
a globally defined pluriharmonic function. Compactness makes this
function constant. In a local flat unitary frame the section is
holomorphic of constant absolute value, hence constant, so it is
parallel. A positive power of $\chi$ is therefore trivial.
This recovers finiteness on this restricted locus without requiring
a common exponent for different loci.
